\documentclass[10pt,a4paper]{amsart}
\usepackage[margin=19mm]{geometry}
\usepackage{amsmath,amssymb,mathtools}
\usepackage[scr=rsfs,cal=euler]{mathalfa}
\usepackage{xfrac}
\usepackage[unicode,hidelinks,bookmarks=false]{hyperref}
\newcommand{\ie}{i.e.\ }
\newcommand{\cf}{cf.\ }
\newcommand{\resp}{respectively\ }
\newcommand{\Subsection}{\subsection}
\providecommand{\nobreakdash}{\nobreak\hbox{-}}
\setlength{\emergencystretch}{2em}
\multlinegap0pt
\usepackage{mathtools}
\usepackage{amssymb}
\usepackage{enumerate}
\usepackage{tikz}
\usepackage{dsfont}
\newtheorem{proposition}{Proposition}
\newcommand{\Bb}{{\mathcal B}}
\newcommand{\CM}{{\mathbb C}}
\newcommand{\Hh}{{\mathcal H}}
\newcommand{\Xx}{{\mathcal X}}
\newcommand{\Zz}{{\mathcal Z}}
\DeclareMathOperator{\diam}{diam}
\title{Kubo Formulas and Bulk-Edge Correspondence for Curved Boundaries}
\author{Aren Martinian, Tom Stoiber}
\date{Source: arXiv:2608.29524v1 (CC BY 4.0)}
\begin{document}
\maketitle

\noindent\emph{Source and licence.} Kubo Formulas and Bulk-Edge Correspondence for Curved Boundaries. Version arXiv:2608.29524v1 (CC BY 4.0), distributed by arXiv under the Creative Commons Attribution 4.0 International licence, http://creativecommons.org/licenses/by/4.0/, as recorded in the arXiv licence metadata for this exact version and shown on the abstract page https://arxiv.org/abs/2608.29524v1. Full licence notice: \texttt{licenses/arxiv-2608.29524-CC-BY-4.0.txt}.\par\smallskip

\noindent\textit{Editorial scope.} Counted unit: the article's proposition prop:morita\_continuous (source0.tex lines 2881--2896). The article's complete original proof of it is delivered with it (source0.tex lines 2898--2977, one \texttt{proof} environment of 80 lines). No mathematics is written, repaired or re-derived: the statement and the proof are the article's own lines, in the article's own notation, and no retained line is substituted or re-flowed.  No further source-local context is needed: every cross-reference of every retained line resolves inside the two delivered spans.  Nothing was omitted from any retained span, so the package carries no omission record.  No retained line contains a citation, so the wrapper carries no bibliography and no bibliography entry.  No retained line contains a figure environment, an image-inclusion command or drawing code, so no image is delivered and the package declares no asset dependency.  Editorial formatting only, in addition: an AMS \texttt{amsart} preamble that restates the article's own declarations and macros used by the retained lines, namely the environments \texttt{proposition} on separate counters, together with the standard packages those lines need.  The retained environments print in this document as Proposition 1 in order of appearance, whereas the article prints the same items with the numbering of its own full document.  The complete native source of the article as distributed in the arXiv e-print for this exact version is bundled unchanged as \texttt{sources/208861/source0.tex}.
\begin{proposition}\label{prop:morita_continuous}
Let $\Xx$ be a proper metric space with bounded geomety. Let $\Hh_\Xx$ be a coarsely faithful geometric module with $\Hh_\lambda\neq \{0\}$ for all $\lambda\in \Lambda$. Choose unit vectors $\xi_\lambda\in\Hh_\lambda$ and set $p:=\bigoplus_{\lambda\in\Lambda} |\xi_\lambda\rangle\langle\xi_\lambda|$.
Then $p\in C_u^*(\Xx,\Hh_\Xx)$ is a full projection and compression induces an
isomorphism
\[ pC_u^*(\Xx,\Hh_\Xx)p\cong C_u^*(\Lambda).\]
Consequently, $C_u^*(\Xx,\Hh_\Xx)$ and $C_u^*(\Lambda)$ are strongly Morita
equivalent.

For a nonempty closed subset $\Zz \subset \Xx$, for $\Lambda_\Zz \coloneqq \{\lambda \in \Lambda: Q_{\lambda} \cap \Zz \neq \varnothing \}$, $p$ acts as a full multiplier projection on $C^*_u(\Zz \subset \Xx, \Hh_\Xx)$ and
\[pC_u^*(\Zz\subset\Xx,\Hh_\Xx)p
 \cong C_u^*(\Lambda_\Zz\subset\Lambda, \ell^2(\Lambda)).\]
The full-corner inclusions thus induce isomorphisms of $K$-theory groups
\[K_*\bigl(C_u^*(\Lambda, \ell^2(\Lambda))\bigr)
 \cong K_*\bigl(C_u^*(\Xx,\Hh_\Xx)\bigr); \qquad  K_*\bigl(C_u^*(\Lambda_\Zz\subset\Lambda)\bigr)
 \cong K_*\bigl(C_u^*(\Zz\subset\Xx,\Hh_\Xx)\bigr).\]
\end{proposition}

\begin{proof}
We abbreviate $A=C_u^*(\Xx,\Hh_\Xx)$ and, for $T\in\Bb(\Hh_\Xx)$, write
\[T_{\lambda\mu}:=\chi_\lambda T\chi_\mu
 \in\Bb(\Hh_\mu,\Hh_\lambda).\]

Clearly, $p$ has finite propagation and is uniformly locally compact, hence it is in $C^*_u(\Xx,\Hh_\Xx)$.

Let $U:\ell^2(\Lambda)\to p\Hh_\Xx$ be the unitary
$U\delta_\lambda=\xi_\lambda$.  On finite-propagation elements of $C^*_u(\Lambda,\ell^2(\Lambda))$ define as a strongly convergent sum
\[\Phi(a):=UaU^* =\sum_{\lambda,\mu\in\Lambda} a_{\lambda\mu}|\xi_\lambda\rangle\langle\xi_\mu|.\]
The propagation changes by at most $2R$ for $R=\sup_{\lambda}\diam(Q_\lambda)$, and every block of $\Phi(a)$ has rank
at most one, hence $\Phi(a)\in A$. Likewise, if $T\in A$ has finite propagation then $pTp$ is in the image of $\Phi$. Taking the closure one obtains an isomorphism $pAp\cong C_u^*(\Lambda,\ell^2(\Lambda)).$

It remains to prove that $p$ is full. 
By bounded geometry, we get for $s\geq0$ that
\[N_s:=\sup_{\lambda\in\Lambda} \#\bigl(B_s(\lambda)\cap\Lambda\bigr)<\infty.\]
 Let $\mathcal A_{\mathrm{fr}}\subset \CM_u(\Xx, \Hh_\Xx)$ be the set of finite-propagation operators $T$ for which $\sup_{\lambda,\mu}\operatorname{rank}(T_{\lambda\mu})<\infty$. 
We claim that $\mathcal A_{\mathrm{fr}}$ is dense in $A$.  It is enough to
consider a finite-propagation uniformly locally compact operator $T$.  Put
$S=\operatorname{prop}(T)$ and $D=N_{S+2R}$.  If $T_{\lambda\mu}\neq0$, then
$d(\lambda,\mu)\leq S+2R$, so every row and every column of the block matrix of
$T$ contains at most $D$ nonzero blocks.

Fix any $\epsilon>0$ and put $\delta=\epsilon/D$. Uniform local compactness, applied to the family
$\{\chi_{B_R(\lambda)}T:\lambda\in\Lambda\}$, gives an integer $n$ and
finite-rank operators $G_\lambda=\chi_\lambda G_\lambda$ of rank at most $n$ such that
\[\|\chi_{\lambda}T-G_\lambda\|<\delta
 \qquad \forall\lambda\in\Lambda.\]
Define a block matrix $F$ by
\[F_{\lambda\mu}:=
 \begin{cases}
  \chi_\lambda G_\lambda\chi_\mu,
     &d(\lambda,\mu)\leq S+2R,\\
  0,&d(\lambda,\mu)>S+2R.
 \end{cases}\]
Then $\operatorname{rank}(F_{\lambda\mu})\leq n$ and
$\|T_{\lambda\mu}-F_{\lambda\mu}\|<\delta$ for every $\lambda,\mu$.
The Schur test gives
\[\|T-F\|\leq D\delta=\epsilon.\]
The operator $F$ is in $\mathcal A_{\mathrm{fr}}$, hence this proves the density claim.

We now show that every $F\in\mathcal A_{\mathrm{fr}}$ is a linear combination of elements in $ApA$.  Let
$n$ bound the ranks of its blocks and let $D$ bound the number of nonzero
blocks in every row and column.  We can think of
\[\mathcal E_F:=\{(\lambda,\mu):F_{\lambda\mu}\neq0\}\]
as the edge set of a bipartite graph with two copies of $\Lambda$.  This graph
has degree at most $D$.  Since it is countable, a greedy edge-colouring with
$2D-1$ colours decomposes $\mathcal E_F$ into finitely many matchings
$\mathcal M_1,\ldots,\mathcal M_{2D-1}$.

For each
$(\lambda,\mu)\in\mathcal E_F$, choose a singular-value decomposition
\[
F_{\lambda\mu} =\sum_{j=1}^n s_{\lambda\mu,j} |\eta_{\lambda\mu,j}\rangle \langle \zeta_{\lambda\mu,j}|,\]
where $\eta_{\lambda\mu,j}\in\Hh_\lambda$ and
$\zeta_{\lambda\mu,j}\in\Hh_\mu$ are unit vectors, and the singular values satisfy
$0\leq s_{\lambda\mu,j}\leq\|F_{\lambda\mu}\|$.  For a matching
$\mathcal M_c$ and $1\leq j\leq n$, set
\begin{align*}
 x_{c,j}:=\sum_{(\lambda,\mu)\in\mathcal M_c}
     s_{\lambda\mu,j}
    |\eta_{\lambda\mu,j}\rangle \langle \xi_\lambda|, \qquad
 y_{c,j}:=\sum_{(\lambda,\mu)\in\mathcal M_c}
     |\xi_\lambda\rangle \langle \zeta_{\lambda\mu,j}|.
\end{align*}
Because $\mathcal M_c$ is a matching, these sums
are orthogonal direct sums; in particular,
\[\|x_{c,j}\| \leq \|F\|, \qquad \|y_{c,j}\| \leq 1.\]
They have finite propagation and are locally rank-one, so
$x_{c,j},y_{c,j}\in A$.  Moreover, $x_{c,j}p=x_{c,j}$ and
$py_{c,j}=y_{c,j}$, and blockwise one has
\begin{equation}
\label{eq:Fdecomp}F=\sum_{c=1}^{2D-1}\sum_{j=1}^n x_{c,j}p y_{c,j}.
\end{equation}
Thus $\mathcal A_{\mathrm{fr}}\subset \operatorname{span}(ApA)$.  Since
$\mathcal A_{\mathrm{fr}}$ is dense in $A$, the ideal generated by $p$ is all
of $A$. Equivalently stated, $p$ is a full projection.

For the relative case, let $I_\Zz=C_u^*(\Zz\subset\Xx,\Hh_\Xx)$. If $F\in\mathcal A_{\mathrm{fr}}$ is supported near $\Zz$, then the factorization \eqref{eq:Fdecomp} is already term-wise in $I_\Zz p I_\Zz$ as all non-zero $x_{c,j}$ or $y_{c,j}$ are supported in some thickening of $\Lambda_\Zz$. Hence $I_\Zz$ is contained in the closed linear span of $I_\Zz p I_\Zz$.
\end{proof}

\end{document}
